\chapter{Introducción}
En el área de Diseño de Algoritmos, los métodos de ordenamiento y búsqueda son herramientas importantes para trabajar con conjuntos de datos de manera eficiente. Sin embargo, su importancia no solo está en conocer su funcionamiento teórico, sino también en observar cómo se comportan al aplicarlos sobre datos reales y en distintos escenarios de ejecución como el peor, el mejor y el caso promedio.

En esta tarea se desarrolló un sistema en lenguaje C para gestionar información de deportistas. El programa permite generar registros de manera automática, almacenarlos en archivos CSV, cargarlos para su procesamiento, ordenarlos según distintos campos y realizar búsquedas por ID, además de incorporar la opción de mostrar rankings de deportistas utilizando el puntaje como criterio principal.

Para resolver el problema, se implementaron los algoritmos de ordenamiento Bubble Sort optimizado, Insertion Sort, Selection Sort con optimización de intercambios innecesarios y Cocktail Shaker Sort, junto con búsqueda secuencial y búsqueda binaria iterativa. A partir de estas implementaciones, también se realizaron pruebas experimentales para comparar sus tiempos de ejecución en distintos tamaños de entrada.

El propósito de este informe es presentar el desarrollo del sistema, describir los algoritmos utilizados y analizar sus resultados teóricos y experimentales. De esta forma, se busca relacionar la implementación práctica con el comportamiento de cada algoritmo, identificando sus diferencias según el caso elegido.